\section{Билет 39. Основные положения молекулярно-кинетической теории (МКТ). Основное уравнение МКТ с выводом. Постоянная Больцмана.}

\begin{definition}
\textbf{Молекулярно-кинетическая теория (МКТ)} --- раздел физики, изучающий строение и свойства веществ на основе представлений о том, что все тела состоят из мельчайших частиц (молекул и атомов), находящихся в непрерывном хаотическом движении.
\par\noindent\textbf{Три основных положения МКТ:}
\begin{enumerate}
    \item Все вещества --- жидкие, твёрдые и газообразные --- образованы из мельчайших частиц: молекул, которые сами состоят из атомов. Молекулы и атомы электрически нейтральны.
    \item Атомы и молекулы находятся в непрерывном хаотическом тепловом движении.
    \item Частицы взаимодействуют друг с другом силами электрической природы; гравитационное взаимодействие между ними пренебрежимо мало.
\end{enumerate}
\end{definition}

\begin{theorem}[Вывод основного уравнения МКТ]
Рассмотрим идеальный газ в сосуде кубической формы с ребром $l$. Пусть в единице объёма содержится $n$ молекул (концентрация), каждая массой $m_0$.

\textbf{1. Импульс, передаваемый одной молекулой.}
Молекула, движущаяся перпендикулярно стенке со скоростью $v_x$, при абсолютно упругом ударе изменяет свою проекцию импульса на $2m_0 v_x$. По закону сохранения импульса, стенка получает импульс:
\begin{equation}
    \Delta p_x = 2 m_0 v_x.
\end{equation}

\textbf{2. Число ударов о стенку за время $\Delta t$.}
За время $\Delta t$ о стенку площадью $S$ успеют удариться только те молекулы, которые находятся в объёме $S \cdot v_x \Delta t$ и движутся к стенке. Учитывая, что движение хаотично, к данной стенке летит $1/6$ всех молекул (по трём осям, в двух направлениях):
\begin{equation}
    N_{\text{уд}} = \frac{1}{6} n S v_x \Delta t.
\end{equation}

\textbf{3. Сила давления.}
Общий импульс, переданный стенке за время $\Delta t$:
\begin{equation}
    \Delta P = N_{\text{уд}} \cdot 2 m_0 v_x = \frac{1}{6} n S v_x \Delta t \cdot 2 m_0 v_x = \frac{1}{3} n m_0 v_x^2 S \Delta t.
\end{equation}
Сила --- это изменение импульса в единицу времени:
\begin{equation}
    F = \frac{\Delta P}{\Delta t} = \frac{1}{3} n m_0 v_x^2 S.
\end{equation}

\textbf{4. Давление.}
Давление --- сила, действующая на единицу площади:
\begin{equation}
    p = \frac{F}{S} = \frac{1}{3} n m_0 \langle v_x^2 \rangle.
\end{equation}
Учитывая, что для хаотического движения $\langle v^2 \rangle = \langle v_x^2 \rangle + \langle v_y^2 \rangle + \langle v_z^2 \rangle = 3\langle v_x^2 \rangle$, получаем \textbf{основное уравнение МКТ}:
\begin{equation}
    p = \frac{1}{3} n m_0 \langle v^2 \rangle = \frac{2}{3} n \left\langle \frac{m_0 v^2}{2} \right\rangle = \frac{2}{3} n \langle E_{\text{пост}} \rangle. \label{eq:mkt_main_39}
\end{equation}
\end{theorem}

\begin{definition}
\textbf{Постоянная Больцмана} $k$ --- фундаментальная физическая константа, связывающая температуру с энергией на молекулярном уровне:
\begin{equation}
    k = \frac{R}{N_A} = \frac{8.31}{6.02 \cdot 10^{23}} \approx 1.38 \cdot 10^{-23} \, \frac{\text{Дж}}{\text{К}}. \label{eq:boltzmann_39}
\end{equation}
где $R$ --- универсальная газовая постоянная, $N_A$ --- число Авогадро.
\end{definition}

\begin{remark}
Основное уравнение МКТ \eqref{eq:mkt_main_39} связывает макроскопический параметр (давление $p$) с микроскопическими характеристиками газа (концентрацией $n$ и средней кинетической энергией поступательного движения молекул $\langle E_{\text{пост}} \rangle$). Это уравнение является фундаментальным для молекулярно-кинетической теории и позволяет объяснить газовые законы с позиций движения молекул.
\end{remark}


\begin{figure}[htbp]
    \centering
    \begin{tikzpicture}[
        scale=1.2,
        >=Stealth,
        font=\small,
        molecule/.style={circle, fill=blue!60, inner sep=2pt},
        wall/.style={fill=gray!30, draw=black, thick}
    ]

    % ==========================================
    % ЧАСТЬ 1: УДАР МОЛЕКУЛЫ О СТЕНКУ
    % ==========================================
    \begin{scope}[shift={(-4, 0)}]
        \node[font=\bfseries, align=center] at (0, 2.2) {Удар молекулы о стенку};

        % Стенка (вертикальная)
        \fill[gray!30] (1.5, -1.2) rectangle (2.0, 1.2);
        \draw[thick] (1.5, -1.2) rectangle (2.0, 1.2);
        \node[right, align=left] at (2.0, 1.0) {стенка\\площадь $S$};

        % Молекула ДО удара (слева)
        \fill[blue!60] (-1.5, 0.5) circle (4pt);
        \draw[->, blue, thick] (-1.4, 0.5) -- (-0.3, 0.5) node[midway, above] {$v_x$};
        \node[below, align=center] at (-1.5, 0.2) {до удара:\\$v_x > 0$};

        % Молекула ПОСЛЕ удара (справа, отлетает)
        \fill[blue!60] (0.72, 0.5) circle (4pt);
        \draw[->, red, thick] (0.6, 0.5) -- (-0.3, 0.5) node[midway, above] {$-v_x$};
        \node[below, align=center] at (0.5, 0.2) {после удара:\\$v_x < 0$};

        % Изменение импульса (стрелка внизу)
        \draw[<->, green!60!black, thick] (-1.0, -0.8) -- (1.0, -0.8)
            node[midway, below, font=\small] {$\Delta p_x = 2m_0 v_x$};

        % Ось x
        \draw[->, thick] (-2.0, -1.5) -- (2.2, -1.5) node[right] {$x$};
        \node at (-2.0, -1.7) {$O$};
    \end{scope}
    \end{tikzpicture}
    \caption{Модель идеального газа в МКТ. Слева: упругий удар молекулы о стенку передаёт импульс $\Delta p_x = 2m_0 v_x$. Справа: кубический сосуд со стороной $l$ содержит $N$ молекул, движущихся хаотично.}
    \label{fig:mkt_derivation}
\end{figure}

\begin{figure}[htbp]
    \centering
    \begin{tikzpicture}[
        scale=1.3,
        >=Stealth,
        font=\small,
        molecule/.style={circle, fill=blue!60, inner sep=2pt}
    ]

    \node[font=\bfseries] at (1.5, 3.5) {Кубический сосуд (рёбра $l$)};

    % === КУБ В ИЗОМЕТРИИ ===
    % Координаты вершин (передняя грань и задняя со смещением)
    % Передняя грань: A(0,0), B(2,0), C(2,2), D(0,2)
    % Задняя грань: E(0.7,0.7), F(2.7,0.7), G(2.7,2.7), H(0.7,2.7)

    % Задняя грань (рисуем первой, чтобы была "сзади")
    \draw[thick, gray!50] (0.7,0.7) -- (2.7,0.7) -- (2.7,2.7) -- (0.7,2.7) -- cycle;

    % Соединительные рёбра
    \draw[thick] (0,0) -- (0.7,0.7);
    \draw[thick] (2,0) -- (2.7,0.7);
    \draw[thick] (2,2) -- (2.7,2.7);
    \draw[thick] (0,2) -- (0.7,2.7);

    % Передняя грань (рисуем последней, чтобы была "спереди")
    \draw[thick, fill=blue!5, fill opacity=0.3] (0,0) -- (2,0) -- (2,2) -- (0,2) -- cycle;
    \draw[thick] (0,0) -- (2,0) -- (2,2) -- (0,2) -- cycle;

    % === ПОДПИСИ РАЗМЕРОВ l ===
    % Нижнее ребро (внизу под кубом)
    \draw[<->] (0, -0.35) -- (2, -0.35) node[midway, below] {$l$};

    % Левое вертикальное ребро (слева от куба)
    \draw[<->] (-0.35, 0) -- (-0.35, 2) node[midway, left] {$l$};

    % Глубина (сверху, вдоль наклонного ребра)
    \draw[<->] (0.35, 2.35) -- (1.05, 3.05) node[midway, above left] {$l$};

    % === МОЛЕКУЛЫ ВНУТРИ ===
    % Молекула 1 (ближе к передней грани)
    \fill[blue!60] (0.6, 0.8) circle (3pt);
    \draw[->, blue, thick] (0.6, 0.8) -- (1.0, 1.0);

    % Молекула 2 (в центре)
    \fill[blue!60] (1.3, 1.2) circle (3pt);
    \draw[->, blue, thick] (1.3, 1.2) -- (1.1, 1.5);

    % Молекула 3 (ближе к задней грани)
    \fill[blue!60] (1.8, 1.6) circle (3pt);
    \draw[->, blue, thick] (1.8, 1.6) -- (2.0, 1.4);

    % Молекула 4 (сверху)
    \fill[blue!60] (1.0, 1.7) circle (3pt);
    \draw[->, blue, thick] (1.0, 1.7) -- (1.3, 1.5);

    % Молекула 5 (справа внизу)
    \fill[blue!60] (1.7, 0.6) circle (3pt);
    \draw[->, blue, thick] (1.7, 0.6) -- (1.5, 0.9);

    % === БЛОК ПАРАМЕТРОВ (справа от куба, с отступом) ===
    \node[
        draw,
        fill=yellow!10,
        rounded corners,
        align=left,
        font=\small,
        inner sep=6pt,
        text width=3.2cm
    ] at (4.2, 1.3)
    {
        \textbf{Параметры сосуда:}\\[6pt]
        Площадь стенки:\\
        \quad $S = l^2$\\[6pt]
        Объём:\\
        \quad $V = l^3$\\[6pt]
        Концентрация:\\
        \quad $n = \dfrac{N}{V}$
    };

    % Стрелка от куба к блоку (опционально, для связи)
    \draw[->, gray, dashed] (2.2, 1.0) -- (2.8, 1.3);

    \end{tikzpicture}
    \caption{Кубический сосуд со стороной $l$, в котором находятся $N$ молекул идеального газа, движущихся хаотично (синие стрелки).}
    \label{fig:cubic_vessel}
\end{figure}

\vspace{0.5cm}

\begin{figure}[htbp]
    \centering
    \begin{tikzpicture}[
        scale=1.1,
        >=Stealth,
        font=\small,
        thick
    ]

    \node[font=\bfseries] at (4, 6.5) {Распределение молекул по скоростям (Максвелл)};

    % === ОСИ ===
    \draw[->] (0, 0) -- (8, 0) node[right, font=\normalsize] {$v$};
    \draw[->] (0, 0) -- (0, 5.5) node[above, font=\normalsize] {$f(v)$};
    \node[below left] at (0, 0) {$O$};

    % === КРИВАЯ РАСПРЕДЕЛЕНИЯ ===
    % f(v) = 2.2 * v^2 * exp(-0.21 * v^2) — пик при v≈2.2, высота ≈3.9
    \draw[blue, very thick, domain=0:7.5, samples=200]
        plot (\x, {2.2 * \x^2 * exp(-0.21 * \x^2)});

    % === ХАРАКТЕРИСТИЧЕСКИЕ СКОРОСТИ ===
    \def\vmp{2.2}    % наиболее вероятная (пик кривой)
    \def\vmean{2.9}  % средняя
    \def\vrms{3.5}   % среднеквадратичная

    % Вертикальные линии
    \draw[red, dashed, thick] (\vmp, 0) -- (\vmp, 3.9);
    \draw[green!60!black, dashed, thick] (\vmean, 0) -- (\vmean, 3.2);
    \draw[purple, dashed, thick] (\vrms, 0) -- (\vrms, 2.1);

    % Подписи под осью
    \node[red, below, align=center, font=\footnotesize] at (\vmp, -0.4)
        {$v_{\text{н.в.}}$\\};

    \node[green!60!black, below, align=center, font=\footnotesize] at (\vmean, -0.4)
        {$\langle v \rangle$\\};

    \node[purple, below, align=center, font=\footnotesize] at (\vrms, -0.4)
        {$v_{\text{кв}}$\\};

    % === ПОЯСНЕНИЕ В ПРАВОМ ВЕРХНЕМ УГЛУ ===
    \node[
        draw,
        fill=white,
        rounded corners,
        align=left,
        font=\small,
        inner sep=5pt
    ] at (6.5, 4.5)
    {
        Хаотическое движение:\\[6pt]
        $\langle v_x^2 \rangle = \langle v_y^2 \rangle = \langle v_z^2 \rangle$\\[6pt]
        $\langle v^2 \rangle = 3\langle v_x^2 \rangle$
    };

    % === ПОЯСНЕНИЕ ПОРЯДКА СКОРОСТЕЙ ===
    \node[
        draw,
        fill=yellow!10,
        rounded corners,
        align=center,
        font=\small,
        inner sep=5pt
    ] at (4, -1.8)
    {
        Порядок скоростей:\\
        $v_{\text{н.в.}} < \langle v \rangle < v_{\text{кв}}$
    };

    \end{tikzpicture}
    \caption{Распределение Максвелла: не все молекулы имеют одинаковую скорость. Показаны наиболее вероятная $v_{\text{н.в.}}$, средняя $\langle v \rangle$ и среднеквадратичная $v_{\text{кв}}$ скорости.}
    \label{fig:maxwell_distribution_fixed}
\end{figure}

